\chapter{Examples}
\label{app:examples}
This appendix collects worked examples that complement the formal development of \emph{FELICS}. They illustrate consequences of the normalization procedure that are easier to inspect in concrete numerical settings than in the general derivations of the main text.

\section*{Effects of Normalization on Pearson Correlation}
\begin{example_box}[label=ex:payout_normalization_correlation]
    \textbf{Demonstration:} Payout normalization can alter the Pearson correlation between \emph{CE} and \emph{EE} vectors.

    \smallskip
    \textbf{Setup:} Consider five concepts evaluated using the original metric of Matton et al.~\cite{mattonWalkTalkMeasuring2024} and using FELICS. The underlying answer-distribution changes are held fixed. To isolate the statistical consequence of normalization, the adjusted values are assumed directly rather than derived from particular joint interventions, reference sets, or payout coefficients. Suppose the two methods produce
    \begin{align}
        \mathbf{CE}_{\mathrm{Matton}} & = \left[1,\,1,\,1,\,1,\,2\right], \\
        \mathbf{CE}_{\mathrm{adj}}    & = \left[1.1,\,0.8,\,0.9,\,0.8,\,2.4\right].
    \end{align}
    Both vectors have the same total causal-effect mass and therefore the same mean:
    \begin{equation}
        \sum_{i=1}^{5}\mathrm{CE}_{\mathrm{Matton},i}
        =\sum_{i=1}^{5}\mathrm{CE}_{\mathrm{adj},i}=6,
        \qquad
        \mu_{\mathbf{CE}_{\mathrm{Matton}}}
        =\mu_{\mathbf{CE}_{\mathrm{adj}}}=1.2.
    \end{equation}

    \smallskip
    \textbf{Explanation-Implied Effects:} Assume the LLM’s explanations yield
    \begin{equation}
        \mathbf{EE} = \left[1.0,\, 0.5,\, 1.0,\, 1.0,\, 2.0\right].
    \end{equation}

    \smallskip
    \textbf{Pearson Correlation Computation:} The explanation-effect vector has mean $\mu_{\mathbf{EE}}=1.1$. For the original metric,
    \begin{align}
        \operatorname{cov}(\mathbf{CE}_{\mathrm{Matton}}, \mathbf{EE}) & = \frac{1}{5}\sum_{i=1}^5 (\mathrm{CE}_{\mathrm{Matton},i} - 1.2)(\mathrm{EE}_i - 1.1) = 0.18, \\
        \operatorname{var}(\mathbf{CE}_{\mathrm{Matton}})              & = \frac{1}{5}\sum_{i=1}^5 (\mathrm{CE}_{\mathrm{Matton},i} - 1.2)^2 = 0.16, \\
        \operatorname{var}(\mathbf{EE})                                 & = \frac{1}{5}\sum_{i=1}^5 (\mathrm{EE}_i - 1.1)^2 = 0.24, \\
        \mathrm{PCC}_{\mathrm{Matton}}                                  & = \frac{0.18}{\sqrt{0.16 \cdot 0.24}} \approx 0.919.
    \end{align}
    For the assumed normalized FELICS vector,
    \begin{align}
        \operatorname{cov}(\mathbf{CE}_{\mathrm{adj}}, \mathbf{EE}) & = \frac{1}{5}\sum_{i=1}^5 (\mathrm{CE}_{\mathrm{adj},i} - 1.2)(\mathrm{EE}_i - 1.1) = 0.28, \\
        \operatorname{var}(\mathbf{CE}_{\mathrm{adj}})              & = \frac{1}{5}\sum_{i=1}^5 (\mathrm{CE}_{\mathrm{adj},i} - 1.2)^2 = 0.372, \\
        \mathrm{PCC}_{\mathrm{adj}}                                 & = \frac{0.28}{\sqrt{0.372 \cdot 0.24}} \approx 0.937.
    \end{align}

    \smallskip
    \textbf{Conclusion:} Although normalization preserves the sum and mean of the CE values in this example, it changes how the values are distributed around that mean. Consequently, both the CE variance and its covariance with $\mathbf{EE}$ change, yielding $\mathrm{PCC}_{\mathrm{adj}}\approx0.937$ instead of $\mathrm{PCC}_{\mathrm{Matton}}\approx0.919$. Equality of the aggregate CE score therefore does not imply equality of the resulting Pearson faithfulness score.
\end{example_box}

\section*{Softmax-Ceiling in FELICS in Practice}
\begin{example_box}[label=ex:softmax_ceiling]
    \textbf{Scenario:} Let the prediction space consist of two outcomes, $\mathcal{Y} = \{y_1, y_2\}$. Suppose the model's baseline internal scores (logits) before any intervention are $\theta_{y_1} = \ln(4)$ and $\theta_{y_2} = \ln(6)$.

    Using the softmax formula (Equation~\ref{eq:culprit_softmax}), the baseline probability for $y_1$ is:
    \begin{equation} P(y_1) = \frac{e^{\ln(4)}}{e^{\ln(4)} + e^{\ln(6)}} = \frac{4}{4 + 6} = 0.40 \quad (40\%) \end{equation}

    \smallskip
    \textbf{Individual Interventions:}
    Assume that intervening on a concept acts additively in the model's logit space.
    \begin{itemize}
        \item \textbf{Concept A:} Increases the logit of $y_1$ by $\Delta_A = \ln(2.25)$.
              \begin{equation} P^{(A)}(y_1) = \frac{4 \cdot e^{\ln(2.25)}}{(4 \cdot e^{\ln(2.25)}) + 6} = \frac{4 \cdot 2.25}{9 + 6} = \frac{9}{15} = 0.60 \quad (60\%) \end{equation}
              \emph{Individual effect of A:} $60\% - 40\% = \mathbf{+20\%}$

        \item \textbf{Concept B:} Independently increases the logit of $y_1$ by $\Delta_B = \ln(3.5)$.
              \begin{equation} P^{(B)}(y_1) = \frac{4 \cdot e^{\ln(3.5)}}{(4 \cdot e^{\ln(3.5)}) + 6} = \frac{4 \cdot 3.5}{14 + 6} = \frac{14}{20} = 0.70 \quad (70\%) \end{equation}
              \emph{Individual effect of B:} $70\% - 40\% = \mathbf{+30\%}$
    \end{itemize}

    \smallskip
    \textbf{Naive Expected Joint Effect:} If probabilities were perfectly additive, intervening on A and B together would shift the probability by $20\% + 30\% = \mathbf{+50\%}$, yielding a final probability of $90\%$.

    \smallskip
    \textbf{Actual Joint Intervention (The Soft Ceiling):}
    When intervening on both concepts, their logit changes combine additively ($\Delta_A + \Delta_B = \ln(2.25) + \ln(3.5) = \ln(7.875)$). The softmax function computes the new joint probability as:
    \begin{equation} P^{(A,B)}(y_1) = \frac{4 \cdot e^{\ln(7.875)}}{(4 \cdot e^{\ln(7.875)}) + 6} = \frac{4 \cdot 7.875}{31.5 + 6} = \frac{31.5}{37.5} = 0.84 \quad (84\%) \end{equation}

    \smallskip
    \emph{Actual joint effect:} $84\% - 40\% = \mathbf{+44\%}$

    \smallskip
    \textbf{Conclusion:} Due to the growing denominator in the softmax function, the actual joint effect ($+44\%$) falls short of the naive additive sum ($+50\%$).
\end{example_box}
